% s-227-p227-m002; original PDF 227, 228
\Aphorism{s-227-p227-m002}{319}
Sowie der Geschmack der heutigen Malerei von den prachtvollen Sonnenuntergängen, den grandiosen Hochgebirgslandschaften zurückkommt und lieber den zarten Offenbarungen einer Herbstbeleuchtung, der stillen langsamen Sprache weiter Ebenen und spärlicher Hügel lauscht, so wünschen wir auch der Poesie, der Menschen- und Seelenschilderung, eine Vertiefung ins Schlichte, heimlich Bedeutsame, das sie bisher zu Gunsten des „interessanten Falls“ vernachlässigt hat. Also keinen Rückfall ins Idiotische, keine Anbetung des Philisters: \Sperr{dieser} Gegensatz zum Pathos und zum Prunkstil ist ebensowenig ernst zu nehmen wie in der Malerei das Stillleben oder die andächtige Wiedergabe eines Gemüsebeetes. Aber ich meine, man solle wieder einmal die Menschenseele da aufsuchen, wo sie scheu, halb und zögernd von sich spricht, nachdem wir soviel in Scene gesetzte, immer auskunftbereite und redefertige Menschlichkeit mit angehört haben. Und ferner wollen wir uns wieder den Katastrophen zuwenden, die sich ohne Krampf und Theaterschrei und Verrenkung aller äusseren Linien abspielen, den unterseelischen Erscheinungen, für die es noch keine Sprache und Rhetorik giebt. Wir wollen die nichtpittoresken Charaktere wiedersehen, die weder im Guten noch im Bösen den Horizont mit Schwüngen und Zacken abschneiden, sondern leise, wie die abendliche Haide, ins Ungewisse verdämmern; wir wollen uns vom decorativen Pomp zur intimen Nüance, von der Wachtparade zur Kammermusik weiterentwickeln. Die grossen Verräther, die monströsen Schuld- und Sühne-Helden, der ganze historische und psychologische Spectakel mit Königen, Priestern, Verschwörern und Volk — alles viel zu laut, zu gewaltsam, zu effectvoll, zu sehr Berner Oberland; unser verwöhntes Auge sucht lieber die Melancholie eines entlegenen Teiches, aus dem wunderliche Nebel aufsteigen und sich auf blassgrüne Wiesen legen. Menschen, die wenig reden, weil sie zuviel erfahren, innerliche Naturen, die jedes Hervorkommen eine schreckliche Selbstüberwindung kostet, sensible Wesen, an denen Töne und Stimmungen lange nachhallen, Meister, die heimlich befehlen, und Gesellen, die stumm und unbewusst gehorchen, dazu das Verhältniss von Mann und Weib, ohne die grellen Pointen der Cultur und doch in hundert Formen verzweigt, ewig räthselhaft und nie dasselbe: oh, wie viel ist noch einem Dichter des \Foreign{pianissimo} zu entdecken vorbehalten!
\par

